\scanpage{5}
\mspage{3)}

dont l'image est le sous-ens

\[
\deleted{\unsure{C_{\mathrm{aff\,st}}}}\quad
\overbrace{C_{\mathrm{aff\,st}}(I_0,\pi_0)}^{\inserted{C_{\mathrm{aff\,st}}^{!}(I_0)}}\subset C_{\mathrm{aff\,st}}(I_0)
\]

\unsure{formé|formée} des $\xi$ tels que la structure polygonale correspondante sur $I_0$, soit la structure polygonale $\pi_0$ donnée sur $I_0$, déduite de la structure polygonale de $S$ par passage au \unsure{quotients|quotient}. Donc on a

\[
\widetilde C_{\mathrm{st}}(S)\longrightarrow C_{\mathrm{aff\,st}}^{!}(I_0)
\]

où le ! indique que $I_0$ est muni de $\pi_0$, i.e. qu'on prend $C_{\mathrm{aff\,st}}(I_0,\pi_0)$. La \deleted{\unsure{type|fibre de}} \inserted{fibre de} $\xi\in C_{\mathrm{aff\,st}}^{!}(I_0)$ s'identifie à l'ens des iso polygonaux

\[
S\mathrel{\overset{\sim}{\rule{2.8em}{0.4pt}}}\widetilde I_0(\xi)
\]

qui \deleted{sont} rendent commutatif le triangle

\diagram{\begin{tikzcd}[ampersand replacement=\&,column sep=small]
 S\arrow[rr]\arrow[dr,"\mathrm{can}"']\&\&\widetilde I_0(\xi)\arrow[dl,"\mathrm{can}"]\\
 \& I_0 \&
\end{tikzcd}}

il est de cardinal $2$ ; \uline{on passe de l'un à l'autre par antipodisme}

\deleted{\uline{Ex} $\operatorname{Card} I_0=5$, $C_{\mathrm{aff\,st}}(I_0)\simeq C_{\deleted{\unsure{\mathrm{aff}}}\mathrm{st}}(I_0\amalg\{i_0\})$ se partage en $6$ classes (orbites sous $\mathfrak S_{I_0}$),}
% archive sheet number: 4
